\chapter{Lipase}

\section*{What Is This Test?}
Lipase is a pancreatic enzyme that catalyzes the hydrolysis of triglycerides into monoglycerides and free fatty acids, which is the rate-limiting and essential step in dietary fat digestion. Unlike amylase, lipase is produced almost exclusively by the pancreatic acinar cells (>99\%), with only tiny amounts from lingual (von Ebner glands), gastric, and intestinal sources. This near-exclusive pancreatic origin gives lipase significantly higher specificity for pancreatic pathology compared to amylase. In acute pancreatitis, lipase leaks from damaged acinar cells into the bloodstream. Serum lipase rises within 4--8 hours of symptom onset (slightly later than amylase), peaks at 24--48 hours, and remains elevated for 8--14 days (compared to amylase's 3--5 days), providing a much wider diagnostic window. Lipase has a serum half-life of approximately 6--14 hours. The slower clearance and prolonged elevation make lipase the preferred test for patients who present >24--48 hours after symptom onset. Current guidelines from the American College of Gastroenterology (ACG), the International Association of Pancreatology, and the American Gastroenterological Association all recommend lipase over amylase as the diagnostic laboratory test of choice for acute pancreatitis \cite{iap2013evidence,tietz2023textbook}. Lipase is also more sensitive (85--100\%) than amylase (70--95\%), particularly in alcoholic pancreatitis and hypertriglyceridemia-induced pancreatitis, where amylase may be falsely normal \cite{yadav2002critical}.

\section*{Alternative Names}
Serum Lipase; Pancreatic Lipase; Blood Lipase; LPS; Triacylglycerol Acylhydrolase

\section*{Principle}
Lipase is measured by enzymatic methods using synthetic triglyceride substrates. The most common method uses 1,2-o-dilauryl-rac-glycero-3-glutaric acid-(6-methylresorufin) ester (DGGR) or other chromogenic diglyceride substrates. Lipase hydrolyzes the substrate, releasing the chromophore methylresorufin, which is intensely colored and measured spectrophotometrically at 580 nm. Colipase and bile acids (deoxycholate or taurodeoxycholate) are added to the reaction mixture to ensure full lipase activity, because lipase requires colipase as a cofactor and bile salts for interfacial activation at the lipid-water interface. Older methods using olive oil emulsion and pH-stat titration (measuring free fatty acids released) or turbidimetric methods are now largely replaced by the chromogenic substrate method. The IFCC has not yet published a standardized reference procedure for lipase, so results from different manufacturers may vary. The laboratory should establish the method-specific reference range (typically 10--60 IU/L, with >3$\times$ ULN considered diagnostic for pancreatitis). The DGGR method is specific for pancreatic lipase. This test is NOT performed by hematology analyzers.

\section*{History}
Lipase was first identified as a fat-splitting enzyme by Claude Bernard in 1856. The clinical use of serum lipase for diagnosing acute pancreatitis was first reported by Cherry and Crandall in 1932. For much of the 20th century, lipase assays were technically difficult and unreliable, making amylase the preferred test. The development of reliable chromogenic substrate methods in the 1980s--1990s transformed lipase measurement, and by the 2000s, lipase had surpassed amylase as the recommended test. The ACG formally recommended lipase over amylase in the 2013 acute pancreatitis guidelines \cite{tenner2013american}.

\section*{Why This Test Is Ordered}
\begin{itemize}
  \item Diagnosis of acute pancreatitis: lipase is the preferred laboratory test. Sensitivity 85--100\%, specificity approximately 95--99\%. A normal lipase essentially excludes acute pancreatitis with >95\% negative predictive value \cite{vissers1999amylase}
  \item Evaluation of acute epigastric abdominal pain with suspected pancreatic etiology (radiation to back, nausea/vomiting, history of gallstones or alcohol use)
  \item Monitoring of pancreatitis resolution: lipase levels trend downward as inflammation resolves, though normalization does not directly correlate with clinical recovery and daily monitoring is not recommended
  \item Evaluation of chronic pancreatitis exacerbations: lipase may be normal in stable chronic pancreatitis but rises during acute flares
  \item Detection of pancreatic injury following ERCP (post-ERCP pancreatitis): lipase checked 2--4 hours post-procedure if abdominal pain develops
  \item Diagnosis of pancreatic trauma (blunt abdominal trauma, especially in children with handlebar injuries) and pancreatic duct obstruction (tumor, stone, stricture)
\end{itemize}

\section*{Primary vs Secondary Test}
Lipase is now the primary diagnostic test for acute pancreatitis, recommended over amylase by ACG, IAP, and AGA guidelines. It is the single best laboratory test for pancreatic evaluation.

\section*{How This Test Is Performed}
A healthcare professional draws blood from a vein into a serum separator tube (gold-top) or lithium heparin tube (green-top). Standard venipuncture technique is used. EDTA or citrate tubes should be avoided as they chelate calcium, which is an essential cofactor for lipase activity, though this effect is minimal if the tube contains excess calcium in the assay reagent. Hemolysis does not significantly affect lipase (unlike amylase). The sample is centrifuged and analyzed on an automated chemistry analyzer. Results are available within 1--2 hours.

\section*{What Preparation Is Needed}
No fasting is required. Inform the provider of: (1) Recent alcohol intake (can mildly elevate lipase). (2) Medications: opiates (morphine may cause mild elevation from sphincter of Oddi spasm, but less so than amylase), corticosteroids, furosemide, thiazide diuretics, valproic acid, azathioprine, and heparin (may release lipase from endothelial binding sites). (3) Recent ERCP. (4) Renal disease: lipase is cleared by the kidney (smaller molecular weight than amylase), so CKD can elevate lipase (typically 1.5--2$\times$ ULN in ESRD). (5) A normal lipase does NOT exclude other serious abdominal pathology (perforated viscus, bowel obstruction, mesenteric ischemia, cholecystitis). Lipase is specific for pancreatitis, not for all causes of acute abdominal pain.

\section*{Contraindications and Precautions}
\begin{itemize}
  \item No absolute contraindications. Critical considerations:
  \item (1) Lipase has superior sensitivity and specificity for acute pancreatitis compared to amylase, and current guidelines recommend lipase as the first-line test. If only one test is ordered, it should be lipase.
  \item (2) Lipase elevation >3$\times$ ULN with characteristic pain diagnoses acute pancreatitis (in context of the Atlanta criteria). However, lipase can be elevated in non-pancreatic conditions: renal failure (decreased clearance), cholecystitis, choledocholithiasis, intestinal obstruction, perforated peptic ulcer, and splenic infarction. Elevations in these conditions are usually mild (<3$\times$ ULN).
  \item (3) In hypertriglyceridemia-induced pancreatitis, lipase may be more reliable than amylase because triglyceride interference with lipase assays is less pronounced than with amylase assays, though both can be affected by turbidity.
  \item (4) Lipase does not differentiate acute pancreatitis from other pancreatic diseases (chronic pancreatitis, pancreatic cancer)---elevation indicates ongoing pancreatic inflammation but does not specify etiology.
  \item (5) Lipase may be normal in chronic pancreatitis with extensive acinar destruction (``burned-out'' pancreas), similar to amylase.
  \item (6) Lipase levels do not predict pancreatitis severity---a very high lipase (>1,000 IU/L) does not necessarily indicate severe pancreatitis, and mild elevation does not exclude severe disease. Severity is assessed clinically (SIRS, organ failure) and by imaging (CT Balthazar score, presence of necrosis). Daily trending of lipase in hospitalized pancreatitis patients is not recommended unless there is clinical deterioration.
  \item (7) Macro-lipase (lipase complexed with immunoglobulin) is a rare cause of persistent hyperlipasemia without clinical pancreatitis. PEG precipitation confirms.
\end{itemize}
\section*{Risks}
Minimal risk from venipuncture. Mild pain or bruising at site resolves in 1--2 days.

\section*{What Happens After the Test}
Results available within 1--2 hours. A lipase >3$\times$ ULN (typically >160--200 IU/L, method-dependent) with acute epigastric pain strongly supports acute pancreatitis. According to the revised Atlanta classification (2012), two of three criteria are required for acute pancreatitis diagnosis: (1) characteristic abdominal pain, (2) serum lipase (or amylase) >3$\times$ ULN, (3) characteristic findings on abdominal imaging (CT, MRI, or US) \cite{banks2013classification}. A normal lipase essentially excludes acute pancreatitis with >95\% NPV. Management: aggressive IV fluid resuscitation (Lactated Ringer's preferred over normal saline, 5--10 mL/kg/hour initially, 250--500 mL/hour for most adults), NPO or early oral feeding (mild pancreatitis---early oral feeding within 24 hours is now recommended over prolonged NPO), pain control (IV opioids), and identification and management of etiology (gallstone removal by ERCP if cholangitis present, cholecystectomy for gallstone pancreatitis, alcohol cessation) \cite{tierney2023current}. Severity scoring (Ranson criteria at 48 hours, BISAP at 24 hours, APACHE II) identifies severe pancreatitis requiring ICU admission.

\section*{Understanding Results}

\subsection*{Below Normal}
\begin{itemize}
  \item Low lipase is not clinically significant. Lipase is normally very low (<10 IU/L in some assays).
  \item Chronic pancreatitis with extensive acinar cell loss (``burned-out'' pancreas) results in low or normal lipase despite significant pancreatic damage.
  \item Cystic fibrosis with pancreatic insufficiency: lipase production is markedly reduced.
  \item A normal lipase essentially excludes acute pancreatitis (>95\% NPV).
\end{itemize}

\subsection*{Above Normal}
\begin{itemize}
  \item \textbf{Acute pancreatitis (>3$\times$ ULN):} Most common causes: gallstones (40\%), alcohol (30\%), idiopathic (15\%), hypertriglyceridemia (>1,000 mg/dL), post-ERCP, medications, trauma, hypercalcemia, autoimmune pancreatitis, pancreatic cancer, viral infections, scorpion stings, genetic (PRSS1, SPINK1, CFTR mutations). Lipase is superior to amylase for: (a) late presentation (elevated for 8--14 days vs. 3--5 days for amylase), (b) alcoholic pancreatitis (amylase production is lower in alcoholics due to chronic acinar damage), (c) hypertriglyceridemic pancreatitis (triglycerides interfere less with lipase assays than amylase).
  \item \textbf{Mild-moderate elevation (1--3$\times$ ULN):} (a) Non-pancreatic abdominal pathology: acute cholecystitis, choledocholithiasis (gallstone in CBD without pancreatitis), intestinal obstruction, perforated peptic ulcer, mesenteric ischemia, acute appendicitis. Pathophysiology: inflammation adjacent to pancreas or intestinal absorption of pancreatic enzymes. Elevations are typically <3$\times$ ULN. (b) Renal failure: decreased renal clearance of lipase. In ESRD, lipase may be 1.5--2$\times$ ULN chronically. A lipase >3$\times$ ULN in a dialysis patient still suggests acute pancreatitis. (c) Diabetic ketoacidosis: transient, mild elevation; resolves with treatment. (d) Medications: opiates (sphincter of Oddi spasm, mild elevation), corticosteroids, furosemide, thiazides, valproic acid, azathioprine, heparin. (e) Post-ERCP: mild transient elevation in 30--70\%; >3$\times$ ULN with pain defines post-ERCP pancreatitis (occurs in 1--10\% of ERCPs). (f) Pancreatic malignancy: mild elevation from duct obstruction or tumor-associated pancreatitis. (g) Abdominal trauma: pancreatic contusion or transection.
  \item \textbf{Macro-lipasemia:} Persistent, asymptomatic elevation of lipase from lipase-immunoglobulin complexes; confirmed by PEG precipitation. Benign---no treatment.
  \item \textbf{Important:} Lipase levels do NOT correlate with pancreatitis severity. A lipase of 10,000 IU/L does not indicate more severe disease than 300 IU/L. Severity is determined by clinical parameters (SIRS, organ failure), Ranson criteria, BISAP score, APACHE II, and CT findings.
\end{itemize}

\section*{Normal Results}
\begin{itemize}
  \item \textbf{Adults:} 10--60 IU/L (varies by method and manufacturer; typical <60--160 IU/L depending on assay) \cite{wallach2022interpretation}
  \item \textbf{Children:} Similar to adult reference range
  \item \textbf{Diagnostic threshold for pancreatitis:} >3$\times$ ULN (typically >160--200 IU/L)
  \item \textbf{Lipase peak in pancreatitis:} Rises 4--8 hours; peaks 24--48 hours; normalizes 8--14 days
  \item \textbf{Lipase half-life:} 6--14 hours
  \item \textbf{Amylase-to-lipase correlation:} Lipase stays elevated 4--5$\times$ longer than amylase
  \item \textbf{Acute pancreatitis diagnosis (Atlanta criteria):} Two of: (1) characteristic pain, (2) lipase/amylase >3$\times$ ULN, (3) imaging findings
\end{itemize}

\section*{Follow-up Tests}
\begin{itemize}
  \item \textbf{Amylase:} Complementary pancreatic enzyme; lipase is preferred but amylase may add diagnostic confidence when both are elevated.
  \item \textbf{Abdominal CT with IV contrast (pancreatic protocol):} Gold standard imaging for pancreatitis diagnosis, severity assessment, and detection of complications (necrosis, pseudocyst, abscess, pseudoaneurysm, splenic vein thrombosis). Best performed 48--72 hours after admission to assess for necrosis (early CT underestimates necrosis).
  \item \textbf{Right upper quadrant ultrasound:} First-line for gallstone pancreatitis; evaluates gallbladder stones, CBD dilation, and choledocholithiasis.
  \item \textbf{MRCP or EUS:} For suspected choledocholithiasis, pancreatic duct abnormalities, or pancreatic mass.
  \item \textbf{Serum triglycerides:} If hypertriglyceridemia-induced pancreatitis suspected (>1,000 mg/dL).
  \item \textbf{Serum calcium:} If hypercalcemia-induced pancreatitis suspected.
  \item \textbf{Liver function tests (ALT, AST, ALP, bilirubin):} Concurrent choledocholithiasis often elevates ALT, ALP, and bilirubin.
\end{itemize}

\section*{References}
\bibliographystyle{unsrtnat}
\bibliography{references}

\clearpage
\section*{Regulatory Status and Authorization Date}
This chapter describes a test category, not a specific commercial product. FDA, CDSCO, EU IVDR, and MHRA approval or registration dates therefore cannot be assigned without the assay manufacturer, product name, instrument, and intended use. For laboratory-developed tests, an individual agency approval date may not exist. See the Preface for the interpretation of regulatory status.

\clearpage
